\chapter{Analysis and history}\label{ch:analysis}

The \ui{Analysis} and \ui{History} tabs, on the right of the step tabs, are for looking back: how the
scheduling of this period went, how the algorithms performed, and what changed in the data. They are
especially useful to explain a decision afterwards and to judge whether the optimiser is worth its time.

\section{How this period's scheduling went}

\shotf[0.95\linewidth]{analysis_summary}{Three questions on whether the tool helps.}{analysis_summary}

\begin{description}[style=nextline, leftmargin=1.2em, font=\bfseries\color{brand}]
  \item[First generation to commit] Minutes between the first \ui{Generate options} and the first commit: is
        the decision quick?
  \item[Commits] How many times the period was committed (``re-committed'' when more than once).
  \item[Hand edits / re-optimisations] How many manual edits and ``optimise the rest'' runs were made: does
        the planner have to edit a lot?
  \item[Optimiser beat the rules?] Whether an optimised schedule beat the dispatching rules under the policy:
        is the optimiser's time worth it?
\end{description}

\section{Run matrix}

\shotf[0.85\linewidth]{analysis_matrix}{Run matrix: one column per ``Generate options'' click, one row per algorithm.}{analysis_matrix}

Each column is one batch (one click on \ui{Generate options}, with its time and input revision); each row is an
algorithm. A cell's colour gives the run status (succeeded, failed, solving, queued, cancelled); a star marks
the recommended schedule and a ring the committed one. Click a cell to open that run.

\section{Every schedule ranked under the selected policy}

\shotf[0.95\linewidth]{analysis_ranking}{All schedules of the period, including the dominated EDD schedule.}{analysis_ranking}

Unlike the Options tab, this table lists \emph{every} schedule, with its rank, score (tier~1 · tier~2), late
orders, shortfall, operating cost, machine shifts, idle time, last lot done, run time and solver status. Each
row says whether it is recommended or ``dominated by'' which recommended option. Click a row for a quick look.

\section{Score by component}

\shotf[0.95\linewidth]{analysis_components}{Each objective term's contribution to the score.}{analysis_components}

Each cell is a term's contribution = weight $\times$ value (the raw value is shown small below). The cells of a
tier add up to that tier's score, so you can see exactly why one schedule ranks above another --- in the
example, EDD loses almost entirely on weighted tardiness.

\section{Algorithm speed and quality}

The summary tiles give the fastest algorithm, which algorithm found the best schedule and after how long, the
number of failed runs, and the median budget.

\shotf[0.85\linewidth]{analysis_speed}{Best score found over time (log scales).}{analysis_speed}

\ui{Best score found over time} draws, for each algorithm's best run, the best score found up to each moment of
its run (a step line that only goes down) against run time on a log scale; the dashed line is the budget.
Dispatching rules build once, so they have no line. A line that is still dropping at the budget suggests a
longer budget (\ui{Thorough}) would help; a line that flattens early suggests it would not.

\shotf[0.95\linewidth]{analysis_quality}{Time and quality by algorithm.}{analysis_quality}

\ui{Time and quality by algorithm} lists schedules produced per runs, time per run, when the best was found,
best and median score across seeds, late orders, and the \textbf{optimality gap}: the distance between the
schedule's score and the bound CP-SAT proved (0\% = optimal). The gap only appears when CP-SAT proved a bound
for the whole problem; otherwise it shows a dash or ``weak bound''.

\section{History}

\shot{history}{The History tab: change log and runs by data version.}{history}

\begin{description}[style=nextline, leftmargin=1.2em, font=\bfseries\color{brand}]
  \item[Change log] Every valid save of the period's data creates a new revision; the log records what changed,
        when and by whom (dataset, products, semi-finished, machines, orders --- filter by kind), and every
        committed schedule.
  \item[Runs by data version] Every run keeps its own copy of the data it used. The panel groups runs by data
        revision (\ui{current} or older) and links to their options, so you can always find the schedules made
        on earlier inputs.
\end{description}
